\begin{summarybox}

	\textbf{\textcolor{CSummary}{一句话核心}}

	\textbf{命题的否定是“条件真且结论假”，否命题是“条件和结论都取否定”，原命题等价于逆否命题而非否命题。}

	\medskip
	\textbf{\textcolor{CSummary}{使用条件}}
	\begin{itemize}[leftmargin=*, itemsep=0.5em, topsep=0.4em]
		\item \textbf{条件 1（命题结构）：} 原命题必须是“若 $p$ 则 $q$”形式，$p$、$q$ 是明确的命题或条件。
		\item \textbf{条件 2（量词判定）：} 若命题隐含全称或存在量词，否定时要相应转换量词。
	\end{itemize}

	\medskip
	\textbf{\textcolor{CSummary}{关键提醒}}
	\begin{itemize}[leftmargin=*, itemsep=0.5em, topsep=0.4em]
		\item \textbf{易错点（概念混淆）：} 勿将否命题当作命题的否定；注意否命题不一定与原命题同真。
		\item \textbf{检查项（量词转换）：} 否定全称命题时不能丢掉量词变化；全称变存在，存在变全称，同时否定结论。
	\end{itemize}
\end{summarybox}
